\subsection{The complete weight-thirteen space at level three}
\label{sec:gl2-weight13}

\begin{proposition}[K17-GL2-MODEL]\label{prop:gl2-weight13}
Let $\chi=\chi_{-3}$, and write $q=\exp(2\pi i\tau)$. Define
\[
 a(\tau)=\sum_{\alpha\in\mathbf Z[\zeta_3]}q^{N(\alpha)},\qquad
 b(\tau)=\frac{\eta(3\tau)^9}{\eta(\tau)^3},\qquad
 D(\tau)=\eta(\tau)^6\eta(3\tau)^6.
\]
Then $S_{13}(\Gamma_0(3),\chi)$ has the basis
\[
 f_0=Da^7,\qquad f_1=Da^4b,\qquad f_2=Dab^2.
\]
It is also an integral basis of the lattice of forms with integral Fourier
coefficients. In this ordered basis the classical Hecke operator $T_7$ is
\[
 A_7=
 \begin{pmatrix}
 -138598&-3312&-138\\
 6438528&159482&4968\\
 -173840256&-3219264&-93886
 \end{pmatrix}.
\]
Consequently its characteristic polynomial is
\begin{align}
 P_7(X)
 &=X^3+73002X^2-10736848500X+248682833875000 \notag\\
 &=(X-40250)^2(X+153502), \label{eq:gl2-p7}
\end{align}
and, for $E_7=1+7^{12}=13841287202$,
\begin{align}
 P_7(E_7)
 &=2651744832098840875904898134016 \notag\\
 &=2^{12}3^{14}\,7\,13^4\,37\,43^2\,47^2\,61\,271^2.
 \label{eq:gl2-p7-evaluation}
\end{align}
\end{proposition}

\begin{proof}
Cl\'ery--van der Geer
\cite[paragraph on $\Gamma_1(3)$, p.~12]{K17CleryVanDerGeer2022}
give
\[
 M(\Gamma_1(3))=\mathbf C[\vartheta,\psi^3],\qquad
 \vartheta=\sum_{\alpha\in\mathbf Z[\zeta_3]}q^{N(\alpha)},\quad
 \psi=\eta(3\tau)^3/\eta(\tau),
\]
where $\vartheta$ and $\psi$ have weight one, and identify the odd-weight
part with $M_{2k+1}(\Gamma_0(3),\chi)$. Thus $a=\vartheta$ and $b=\psi^3$
have weights one and three, with exactly the normalizations used above.
Their Remark~15.2, p.~33, identifies $D$ as a weight-six cusp form on
$\Gamma_0(3)$ with trivial character.

We check division by $D$ at every point. The eta function has no zeros
in the upper half-plane, so neither does $D$. Its product expansion starts
with $q$, giving order one at infinity. Moreover
\[
 D\!\left(-\frac1{3\tau}\right)
 =(-i3\tau)^3(-i\tau)^3D(\tau)
 =-27\tau^6D(\tau).
\]
For the determinant-normalized slash operator this is
$D|_6W_3=-D$, where
$W_3=\left(\begin{smallmatrix}0&-1\\3&0\end{smallmatrix}\right)$.
Thus $D$ also has order one at zero. These are the two cusps of
$\Gamma_0(3)$, of widths one and three. Their stabilizer translations
have lower-right entry one, so the character $\chi$ gives integral local
orders at each cusp. It follows that division by $D$ induces an
isomorphism
\[
 S_{13}(\Gamma_0(3),\chi)\ \xrightarrow{\ \sim\ }\
 M_7(\Gamma_0(3),\chi).
\]
The weighted-degree-seven monomials in the polynomial ring
$\mathbf C[a,b]$ are exactly $a^7,a^4b,ab^2$. This proves the claimed
basis and its completeness.

For clarity, the precise integer expansions used in the calculation are
\begin{align}
 a&=1+6\sum_{n\ge1}\left(\sum_{d\mid n}\chi(d)\right)q^n,
 \label{eq:gl2-theta-expansion}\\
 b&=q\prod_{n\ge1}\frac{(1-q^{3n})^9}{(1-q^n)^3},\qquad
 D=q\prod_{n\ge1}(1-q^n)^6(1-q^{3n})^6.
 \label{eq:gl2-eta-expansion}
\end{align}
The theta formula follows directly from unique factorization in the
Eisenstein integers: there are six generators of each nonzero principal
ideal, and the number of ideals of norm $n$ is
$\sum_{d\mid n}\chi(d)$, as follows prime by prime from splitting,
inertia, and ramification at three. The eta expansions follow from their
defining products. Every coefficient is integral. Their leading terms
give $f_0=q+O(q^2)$, $f_1=q^2+O(q^3)$, and
$f_2=q^3+O(q^4)$. More precisely, the first-three-coefficient matrix is
\[
 L=\bigl([q^i]f_j\bigr)_{\substack{1\le i\le3\\0\le j\le2}}
   =\begin{pmatrix}1&0&0\\36&1&0\\513&21&1\end{pmatrix}.
\]
Since $\det L=1$, these three coefficients determine any form in the
space and recover integral basis coordinates whenever they are integers.
This proves the assertion about the integral Fourier-expansion lattice.

For a prime $p\nmid3$, the classical Hecke normalization is
\[
 [q^n](T_pf)=[q^{pn}]f+\chi(p)p^{12}[q^{n/p}]f,
 \tag{*}
\]
where the last coefficient is zero unless $p\mid n$;
see \cite[\S8.2.2, p.~66]{K17RibetStein2011}.
For $p=7$ and $n=1,2,3$, the second term is zero.
The finite products in \eqref{eq:gl2-theta-expansion}--\eqref{eq:gl2-eta-expansion},
calculated modulo $q^{22}$, give
\[
 B_7=\bigl([q^{7i}]f_j\bigr)_{\substack{1\le i\le3\\0\le j\le2}}
 =\begin{pmatrix}
 -138598&-3312&-138\\
 1449000&40250&0\\
 -109731942&-1569198&-60352
 \end{pmatrix}.
\]
Therefore $LA_7=B_7$, proving the displayed matrix.
One can perform the solve without division: for a column $(v_1,v_2,v_3)^t$
of $B_7$, its coordinates are
\[
 c_0=v_1,\qquad c_1=v_2-36c_0,\qquad
 c_2=v_3-513c_0-21c_1.
\]
The trace, the sum of the principal two-by-two minors, and the determinant
of $A_7$ now give \eqref{eq:gl2-p7} by the identity
\[
 \det(XI-A)=X^3-\operatorname{tr}(A)X^2+
 \left(\sum_{i<j}(A_{ii}A_{jj}-A_{ij}A_{ji})\right)X-\det(A).
\]
Substitution of $E_7$ and multiplication give
\eqref{eq:gl2-p7-evaluation}; in particular this integer is nonzero.
Every displayed factor is prime, as is also checked by trial division in
the standalone certificate.

The certificate
\texttt{gl2\_weight13\_certificate.py} reconstructs the products using only
integer arithmetic in the Python standard library. Positive eta exponents
use the finite binomial formula, and negative exponents use
$(1-X)^{-r}=\sum_{k\ge0}\binom{r+k-1}{k}X^k$.
It truncates all operations modulo $q^{56}$, computes $L,B_7,A_7$,
and checks the polynomial and complete integer factorization.
The proof of the basis above makes the first three coefficients
sufficient to identify the operator; the further comparisons through
the coefficient of $q^7$ are additional exact checks, including the
nonzero second term of $(*)$ at $n=7$.
\end{proof}

\paragraph{A second split-prime check.}
The same certificate calculates $T_{13}$ on the entire same space:
\[
 A_{13}=
 \begin{pmatrix}
 -8575918&-182592&-7608\\
 354958848&7857362&273888\\
 -9583888896&-177479424&-6110926
 \end{pmatrix}.
\]
The coefficient matrix used in $LA_{13}=B_{13}$ is
\[
 B_{13}=
 \begin{pmatrix}
 -8575918&-182592&-7608\\
 46225800&1284050&0\\
 -6529199022&-106144518&-4262182
 \end{pmatrix}.
\]
Its characteristic polynomial and evaluation are
\begin{align*}
 P_{13}(X)
 &=X^3+6829482X^2-22485145931700X+15494586622814755000\\
 &=(X-1284050)^2(X+9397582),\\
 P_{13}(1+13^{12})
 &=12646222259779542583967104682915874471936\\
 &=2^{18}3^{14}7^2\,13^4\,19\,37\,47^2\,193\,337^2\,14551^2.
\end{align*}
These identities follow by the same first-three-coefficient solve and
determinant formula, with every required source coefficient included in
the exact truncation.

\paragraph{Arithmetic scope.}
This proposition computes all classical Hecke eigenvalues in the
specified space, including multiplicities. It does not by itself prove
that a Galois-theoretic branch has conductor dividing three, that its
nebentypus is $\chi_{-3}$, or that its trace equals this classical
normalization. Those are separate source-interface assertions. Whenever
that interface places a weight-thirteen eigenform in this space,
its classical $T_7$ eigenvalue is a root of $P_7$.
No conclusion about an algebraic $K$-group is inferred here merely from
the existence or factorization of this polynomial.
